\chapter{2つ目の出力}
% full version: chapters/16_chemoutput.tex

マシンには出力が2つある。速い出力は筋肉だ。遅い出力は体の化学である。脳は心臓、血管、腺を制御し、血液を通して全身に一度にメッセージを送る。血液はあらゆるものの中で最もブロードキャスト的なバスだ。1つのメッセージが約1分であらゆる細胞に届き、その意味は、いつものように受け手の「錠」が与える。

\section*{アクセルとブレーキ}

心臓は自分で拍動する。自前のメトロノームをもっているのだ。脳は、符号が逆の2本の配線でテンポを調整するだけである。\nt{NORA}はアクセルペダルで、速めるがゆっくりだ。\nt{ACET}はブレーキで、すばやく遅くする。長い化学反応の連鎖を必要としないからだ。そしてここでようやく\nt{ADRE}型に仕事が見つかる。血液中に放出されたアドレナリンは同じアクセルペダルだが、全身に対して一度に踏まれる。

\pencilfig{16}{\pencilart{16}}{心臓は符号が逆の2本の配線で操られる。アクセルとブレーキである。}

\section*{フィードバックつきのカスケード}

多くのホルモンはカスケードで制御されている。脳が1つ目の腺に命じ、それが2つ目の腺に命じ、2つ目がホルモンを分泌し、ホルモンが脳に「もう十分だ」と伝える。3段のサーモスタットである。エンジニアはこうした回路の危険を知っている。ループを「敏感」にしすぎると揺れ始め、レベルは保たれずに振動する。だからこの調節器の精度には限界がある。安定でありながら非常に正確にすることはできないのだ。

\pencilfig{127}{%
% three identical first-order stages with proportional feedback: secretion = max(0, K (target - level)).
% K = 3 settles at 3/4 of the target; K = 12 (above the stability limit K = 8) keeps oscillating. x = 0.42 t, y = 0.55 level
\foreach \yo/\t in {1.75/{弱いループ：穏やかだが目標に届かない}, 0/{強いループ：目標に近いが揺れ続ける}}{
  \draw[pen] (0,\yo) -- (8.5,\yo);
  \draw[pcGray, densely dashed, line width=0.5pt] (0,\yo+0.55) -- (8.5,\yo+0.55);
  \node[lbl, anchor=south west] at (2.0,\yo+0.8) {\t};}
\node[lbl, anchor=west] at (8.55,2.3) {目標}; \node[lbl, anchor=west] at (8.55,0.55) {目標};
\begin{scope}[yshift=1.75cm]
\draw[penlite=pcBlue] plot[smooth] coordinates {(0.00,0.00) (0.17,0.01) (0.34,0.08) (0.50,0.19) (0.67,0.33) (0.84,0.46) (1.01,0.55) (1.18,0.59) (1.34,0.58) (1.51,0.55) (1.68,0.49) (1.85,0.43) (2.02,0.37) (2.18,0.34) (2.35,0.33) (2.52,0.34) (2.69,0.37) (2.86,0.40) (3.02,0.43) (3.19,0.44) (3.36,0.45) (3.53,0.45) (3.70,0.44) (3.86,0.42) (4.03,0.41) (4.20,0.40) (4.37,0.39) (4.54,0.39) (4.70,0.40) (4.87,0.40) (5.04,0.41) (5.38,0.42) (5.88,0.42) (6.38,0.41) (7.06,0.41) (8.40,0.41)};
\end{scope}
\draw[penlite=pcRed] plot[smooth] coordinates {(0.00,0.00) (0.17,0.05) (0.34,0.30) (0.50,0.68) (0.67,1.02) (0.84,1.23) (1.01,1.30) (1.18,1.26) (1.34,1.16) (1.51,1.02) (1.68,0.87) (1.85,0.72) (2.02,0.58) (2.18,0.47) (2.35,0.39) (2.52,0.38) (2.69,0.46) (2.86,0.57) (3.02,0.67) (3.19,0.71) (3.36,0.70) (3.53,0.65) (3.70,0.58) (3.86,0.50) (4.03,0.43) (4.20,0.41) (4.37,0.45) (4.54,0.53) (4.70,0.61) (4.87,0.65) (5.04,0.64) (5.21,0.60) (5.38,0.54) (5.54,0.47) (5.71,0.42) (5.88,0.42) (6.05,0.48) (6.22,0.56) (6.38,0.62) (6.55,0.64) (6.72,0.62) (6.89,0.56) (7.06,0.50) (7.22,0.44) (7.39,0.42) (7.56,0.45) (7.73,0.53) (7.90,0.60) (8.06,0.63) (8.23,0.62) (8.40,0.58)};
\node[lbl, anchor=north] at (4.2,-0.05) {時間};
}{3段のホルモン・カスケードのモデル。弱いフィードバックはレベルを落ち着かせるが、目標よりかなり低いままにする。強いフィードバックは目標に近づけるが、レベルはもう保たれず揺れ動く。}

\section*{レベルではなくパルス}

パルスでしか働かないホルモンもある。絶えず与えると受け手の腺は「疲れて」停止してしまうが、1時間に1回の短いパルスで与えると働く。受け手が読んでいるのはレベルではなくリズムだ。モールス符号の章の、変化だけを聞くシナプスと同じである。

\section*{1日}

最後に、体には周期が約1日のゆっくりした発振器がある。その自然な周期は24時間とわずかに異なり、毎朝、光がそれを調整する。ラジオ受信機が局に同調するように。コンピュータのクロック発生器と混同してはいけない。これは計算のステップを刻むのではない。マシン全体の動作モード、つまり昼と夜、覚醒と睡眠を切り替えるのである。

\fullversion{16}
